
% Springer LLNCS paper adapted from MDPI medical text summarization study
% Template: llncs.cls (Lecture Notes in Computer Science)
\documentclass[runningheads]{llncs}
%
\usepackage[T1]{fontenc}
\usepackage{graphicx}
\usepackage{booktabs}
\usepackage{multirow}
\usepackage{amsmath}
\usepackage{hyperref}
\usepackage{xcolor}
\renewcommand\UrlFont{\color{blue}\rmfamily}
\urlstyle{rm}
%
\begin{document}
%
\title{Evaluating Medical Text Summaries Using Automatic Metrics and LLM-as-a-Judge: A Pilot Study}
%
\titlerunning{Evaluating Medical Text Summaries with LLM-as-a-Judge}
% If the paper title is too long for the running head, you can set
% an abbreviated paper title here
%
\author{Yuriy Vasilev\inst{1}\orcidID{0000-0000-0000-0001} \and
Irina Raznitsyna\inst{1}\orcidID{0000-0000-0000-0002} \and
Anastasia Pamova\inst{1,2}\orcidID{0000-0000-0000-0003} \and
Tikhon Burtsev\inst{1} \and
Tatiana Bobrovskaya\inst{1} \and
Pavel Kosov\inst{1} \and
Anton Vladzymyrskyy\inst{1,3} \and
Olga Omelyanskaya\inst{1,2} \and
Kirill Arzamasov\inst{1}}
%
\authorrunning{Vasilev et al.}
% First names are abbreviated in the running head.
% If there are more than two authors, 'et al.' is used.
%
\institute{Research and Practical Clinical Center for Diagnostics and Telemedicine Technologies of the Moscow Health Care Department, 127051 Moscow, Russia \\
\email{y.vasilev@telemed.center} \and
Institute of Artificial Intelligence, MIREA---Russian Technological University, 119454 Moscow, Russia \\
\email{a.pamova@mirea.ru} \and
Department of Information Technology and Medical Data Processing, Sechenov University, 119991 Moscow, Russia}
%
\maketitle              % typeset the header of the contribution
%
\begin{abstract}
Electronic health records (EHRs) remain a vital source of clinical information, yet processing these heterogeneous data is extremely labor-intensive. Summarization of these data using large language models (LLMs) is considered a promising tool to support practicing physicians. Unbiased, automated quality control is crucial for integrating these tools into routine practice, saving time and labor. This pilot study aimed to assess the potential and constraints of self-contained evaluation of summarization quality (without expert involvement) based on automatic evaluation metrics and LLM-as-a-judge. The summaries of text data from 30 EHRs were generated by six open-source low-parameter LLMs. The medical summaries were evaluated using standard automatic metrics (BLEU, ROUGE, METEOR, BERTScore) as well as the LLM-as-a-judge approach using the following criteria: relevance, completeness, redundancy, coherence and structure, grammar and terminology, and hallucinations. Expert evaluation was conducted using the same criteria. The results showed that LLMs hold great promise for summarizing medical data. Nevertheless, neither the evaluation metrics nor LLM judges are reliable in detecting factual errors and semantic distortions (hallucinations). In terms of relevance, the Pearson correlation between the summary quality score and the expert opinions was 0.688. Completely automating the evaluation of medical summaries remains challenging. Further research should focus on dedicated methods for detecting hallucinations, along with investigating larger or specialized models trained on medical texts.

\keywords{Large language model \and Text summarization \and Electronic health records \and LLM-as-a-judge \and Medical informatics}
\end{abstract}
%
%
\section{Introduction}\label{sec:intro}

Using medical history as the primary source of information about patient health is an important step toward accurate diagnosis. Patient interviews do not always deliver comprehensive clinical presentation, and analyzing medical records is labor-intensive and time-consuming.

Large language models (LLMs) are a type of transformer neural networks that have been pre-trained on vast amounts of textual data. LLMs utilize statistical analysis of text and its components as data units or tokens (sequences of words, syllables, and letters). The models predict the most probable continuation for a given sequence of prompt tokens, taking into account syntactic, semantic, and ontological relationships~\cite{roustan2025}. LLMs are capable of processing large volumes of text and generating human-like summaries and interpretations~\cite{vanveen2024}.

LLMs have been drawing increasing attention in medicine as tools for extracting clinically significant information, analyzing and interpreting data, etc.~\cite{nazi2024,vasilev2025scoping}. Research confirms the effectiveness of LLMs as clinical decision support tools in specialized fields such as oncology~\cite{kaiser2025}, radiology~\cite{nakaura2024}, otolaryngology~\cite{filali2025}, and pediatric nephrology~\cite{niel2025}. LLMs facilitate the automation of documentation workflows by summarizing physician consultations based on transcriptions of clinical conversation~\cite{lee2023}. They are increasingly integrated into care planning, medical education, and other organizational processes~\cite{liu2025}.

Today, the integration of AI-powered technologies in medicine is accompanied by understandable mistrust and apprehension. Any LLM poses the risks of illusory conclusions (``hallucinations''), while both input data and evaluation outputs may be biased due to algorithmic limitations and lack of diversity in training datasets. These shortcomings negatively impact the quality of medical summaries, leading to misrepresentation of the patient's condition and ultimately affecting treatment decisions.

Quality control of LLM performance is critical, as the summaries must comply with established requirements. In a number of studies, the quality of LLM summaries is assessed by both domain specialists and established automatic metrics that evaluate the performance of natural language processing: ROUGE~\cite{lin2004}, BLEU~\cite{papineni2002}, METEOR~\cite{lavie2007}, BERTScore~\cite{zhang2019}, etc. Autonomous summary assessment using automatic evaluation metrics alone is impossible, as they only reflect the similarity between two texts: an LLM output and a reference. The automatic evaluation metrics do not allow assessment of coherence or readability and are unable to determine if key information is missing; therefore, they lack correlation with expert assessment~\cite{vanveen2024,tang2023}.

In searching for an autonomous method for summary evaluation, research has suggested using LLMs as evaluators (LLM-as-a-judge)~\cite{zheng2023}. LLM-as-a-judge utilizes user-defined criteria (e.g., coherence, consistency, relevance, accuracy) defined in prompts. One advantage of this approach is that it does not require reference summaries. However, only a few studies report high consistency between LLM-as-a-judge and expert outputs~\cite{chen2024,croxford2025}. Inherent LLM biases also impact the performance of LLM-as-a-judge, preventing them from being considered reliable, autonomous evaluation tools, especially in the highly sensitive medical field~\cite{vanschaik2024,gu2024}.

The primary objective of this work is to determine whether existing approaches to the automated evaluation of generated texts can be adapted to the task of medical text summarization under constraints of limited computational resources and budget. This paper presents the results of a pilot study assessing the quality of medical text summarization using a set of automatic evaluation metrics and open-source LLM-as-a-judge.

\section{Materials and Methods}\label{sec:methods}

This study assessed the summaries of 30 text documents containing electronic health record (EHR) data. The summaries were generated to extract data relevant to radiological reports. We used six state-of-the-art generative models, varying in size (number of parameters), quantization format, and reasoning mode. The reference summaries were created by clinicians and validated by a radiologist. The summaries generated by the LLMs were assessed using automatic evaluation metrics, two LLM-as-a-judges, and experts.

Thus, the pilot study investigated 180 summaries using each of the criteria below. This sample size allowed us to draw preliminary conclusions regarding the validity of the automatic evaluation metrics and the scores assigned by the LLM judges.

All experiments were run on a dual RTX 3090 Ti setup (24 GB VRAM per GPU) (NVIDIA Corp., Santa Clara, CA, USA), which offered sufficient performance for inference of the LLMs.

\subsection{Dataset}\label{subsec:dataset}

The dataset consisted of 30 text files (from 1450 to 4216 words) generated from EHR data. The inclusion criteria for EHR were the availability of the following information:
\begin{itemize}
    \item Two radiology reports for the same anatomical region, acquired with the same modality;
    \item Examination report created by the specialist who referred the patient for the last of these imaging studies;
    \item A report for an additional imaging study acquired using a different modality or on a different anatomical area;
    \item Two values of a lab test (initial and follow-up);
    \item An additional lab test value different from the above;
    \item An examination report by a physician of any specialty not related to imaging studies;
    \item A hospital discharge summary.
\end{itemize}
Only EHRs belonging to patients over 18 years of age were included in the study.

\subsection{LLMs for Summarizers and LLMs-as-a-Judge}\label{subsec:llms}

Models hosted on the Hugging Face platform were used to summarize the EHR data. LLM inclusion criteria were: released no earlier than Q1 2024; available under the Apache 2.0 open-source license; transformer architecture; capable of handling large volumes of text data. The resulting six LLMs are presented in Table~\ref{tab:llms}.

\begin{table}[t]
\caption{The six LLMs utilized for summarization in this study.}\label{tab:llms}
\centering
\begin{tabular}{@{}lll@{}}
\toprule
\textbf{LLM} & \textbf{Model Size} & \textbf{Developer} \\
\midrule
Gemma-2-27b-it & 27 billion & Google LLC \\
Llama-3.3-70B-Instruct\_Q4 & 70 billion & Meta Platforms Inc. \\
Llama-3.1-8b & 8 billion & Meta Platforms Inc. \\
Llama-3.2-3B-Instruct & 3 billion & Meta Platforms Inc. \\
Gemma-3-12b-it & 12 billion & Google LLC \\
Qwen3-32B\_Q4-nt & 32 billion & Qwen Ltd. \\
\bottomrule
\end{tabular}
\end{table}

The Q4 label denotes the use of the level 4 quantization method q4\_K\_M, which reduces the size of neural network weights by reducing the bit depth of floating-point numbers. The nt label denotes the use of the ``no\_think'' mode, which disables the LLM's analytical capabilities. All models were run at zero temperature, which minimizes generation stochasticity.

Two models from different families were selected to act as LLM-as-a-judge: Qwen3-14b and Mistral-small-24b. The first model was chosen to test the hypothesis about estimate bias caused by testing the model on its own data. No specialized fine-tuning of the models was performed.

\subsection{Reference Summary and Expert Review}\label{subsec:expert}

Four expert physicians (three clinicians and one radiologist) with over three years of experience provided reference summaries for 30 texts. The task was designed to extract the data that radiologists could use in abdominal computed tomography (CT) interpretation, namely complaints underlying the referral, medical and life history (comorbidities, bad habits, family history, surgeries), as well as data from lab tests and imaging studies.

Eighteen experts evaluated the summarization performance. Three experts and a radiologist assessed each generated summary according to the following criteria: relevance, recall, redundancy, coherence and structure, grammar and terminology, and hallucinations. The experts used a binary scale (present/absent) to assess the ``Hallucinations'' criterion. The resulting binary values were converted into quantitative scores: 1 point was assigned to summaries containing a hallucination, and 5 points were assigned if no hallucination was present. The remaining criteria were assessed on a five-point Likert scale, where 1 corresponded to the worst result and 5 to the best.

\subsection{Automatic Evaluation Metrics}\label{subsec:metrics}

The following metrics were analyzed:
\begin{enumerate}
    \item \textbf{ROUGE} (Recall-Oriented Understudy for Gisting Evaluation). This class of metrics evaluates the recall and accuracy of summarization on a scale from 0 to 1 by analyzing the match of n-grams and their sequences with a reference summary. ROUGE-1 evaluates the match of unigrams; ROUGE-2 evaluates bigrams; ROUGE-L focuses on the longest common subsequences.
    \item \textbf{BLEU} (Bilingual Evaluation Understudy) uses a weighted combination of n-gram matches to estimate similarity and introduces a brevity penalty.
    \item \textbf{METEOR} (Metric for Evaluation of Translation with Explicit Ordering) accounts for the variability of words' morphological features through stemming and lemmatization.
    \item \textbf{BERTScore} evaluates the similarity between generated and reference summaries by calculating the semantic distance between individual word vectors using a pre-trained BERT model.
\end{enumerate}

\subsection{Statistical Analysis}\label{subsec:stats}

Statistical analysis was performed using Python 3.12 software with the SciPy and pandas libraries and JASP (v.~0.19.3.0). All evaluations considered statistical significance $\alpha = 0.05$. We analyzed the mean ($M$), the standard deviation ($SD$), the median ($Me$), and the interquartile range $[Q1; Q3]$. The normal distribution of quantitative variables was tested using the Shapiro--Wilk test. The significance of differences was determined using the Wilcoxon signed-rank test. The strength and direction of the relationship between two quantitative variables were assessed using the Spearman correlation coefficient ($r_s$). We used multiple linear regression analysis to test the hypothesis regarding the predictability of expert scores.

\section{Results}\label{sec:results}

The average summary generation time was 25~s. In contrast, the automated evaluation by Qwen3-14b took, on average, 200~s, and that by Mistral-small-24b took 85~s. Table~\ref{tab:expert} presents the expert scores for 30 summaries for each LLM.

\begin{table}[t]
\caption{Expert scores of the summary quality for each LLM ($M \pm SD$).}\label{tab:expert}
\centering
\small
\begin{tabular}{@{}lcccccc@{}}
\toprule
\textbf{Criterion} & \textbf{Gemma-2} & \textbf{Llama-3.3} & \textbf{Llama-3.1} & \textbf{Llama-3.2} & \textbf{Gemma-3} & \textbf{Qwen3} \\
 & \textbf{27b} & \textbf{70B} & \textbf{8b} & \textbf{3B} & \textbf{12b} & \textbf{32B} \\
\midrule
Relevance & $2.6 \pm 1.3$ & $3.8 \pm 1.0$ & $3.3 \pm 0.9$ & $2.4 \pm 1.0$ & $3.6 \pm 1.2$ & $3.8 \pm 1.3$ \\
Recall & $2.6 \pm 1.3$ & $3.7 \pm 1.0$ & $3.1 \pm 1.1$ & $2.3 \pm 1.1$ & $3.9 \pm 1.1$ & $3.9 \pm 1.2$ \\
Redundancy & $3.2 \pm 1.4$ & $4.2 \pm 0.8$ & $3.5 \pm 1.1$ & $3.2 \pm 1.2$ & $4.0 \pm 1.0$ & $3.2 \pm 1.3$ \\
Coherence & $3.6 \pm 1.2$ & $4.4 \pm 1.0$ & $3.6 \pm 1.2$ & $3.1 \pm 1.1$ & $4.2 \pm 1.0$ & $4.2 \pm 1.2$ \\
Grammar & $4.5 \pm 0.6$ & $4.7 \pm 0.7$ & $4.4 \pm 0.8$ & $3.6 \pm 1.2$ & $4.7 \pm 0.5$ & $4.3 \pm 1.1$ \\
Hallucinations & $3.8 \pm 1.9$ & $4.6 \pm 1.2$ & $3.6 \pm 1.9$ & $3.4 \pm 2.0$ & $3.9 \pm 1.8$ & $3.8 \pm 1.9$ \\
\bottomrule
\end{tabular}
\end{table}

The expert scores illustrated the LLM capabilities in summarizing EHR data. The Llama family models with 70 billion and 8 billion parameters demonstrated consistently high summarization performance across all evaluation criteria. Interestingly, the Gemma model (27 billion) demonstrated lower overall performance on average compared to the Llama model (8 billion), with the exception of the ``Grammar and Terminology'' criterion.

For practical purposes, agreement between the absolute values of metrics from LLMs and experts can be neglected. The validity criterion for LLM-as-a-judge is the adequate summary ranking. Therefore, to select an LLM judge, we analyzed the correlation of metrics obtained by experts and the LLM judges using Spearman's rank correlation coefficient.

We also tested the hypothesis that the self-assessing model generates overestimated metric values. Qwen3-14b significantly overestimated the relevance, completeness, and redundancy scores compared to Mistral-small-24b. However, both LLMs identified only two of the nine hallucinations detected by the experts. Therefore, the Qwen3-14b LLM metrics were excluded from further analysis due to lower correlation with the expert assessment compared to Mistral-small-24b, in addition to the bias manifested as score overestimation for the models of its class.

\subsection{Performance of the Automatic Evaluation Metrics}\label{subsec:autometrics}

The expectedly low correlation between the automatic evaluation metrics and the expert scores provides evidence of their infeasibility as reliable independent evaluators. These results were not unexpected: evaluation metrics are only capable of assessing similarity to a reference summary, one of many that meet the criteria. Moreover, the BERTScore metric demonstrated a stronger correlation with the expert scores, as it takes into account possible synonymous substitutions and syntactic paraphrases, demonstrating robustness to lexical diversity.

The results demonstrated low correlation between the automatic evaluation metrics and the LLM judges, which indicates a weak relationship. This led us to conclude that each group of metrics evaluates the summary quality from different perspectives. Therefore, a standalone assessment requires using both metric groups, which do not duplicate but rather complement each other, providing a comprehensive summary quality score.

\subsection{Standalone Integrated Assessment of Summarization Performance}\label{subsec:integrated}

To develop an integrated score, we relied on multiple linear regressions using the elimination method to remove the least significant predictors from a model that initially included all possible predictors. The Mistral-small-24b model was selected as the LLM judge.

Table~\ref{tab:regression} presents the three most significant predictors for each criterion score. Multiple correlation coefficients ($R$) were calculated.

\begin{table}[t]
\caption{Multiple correlation coefficients $R$ of various models with expert score metrics. The $R^2$ coefficient of determination and the root mean square error (RMSE) for the full model are shown.}\label{tab:regression}
\centering
\small
\begin{tabular}{@{}lp{3.5cm}ccccc@{}}
\toprule
\textbf{Expert Score} & \textbf{Top Predictors} & \textbf{$R$ (3)} & \textbf{$R$ (BERT)} & \textbf{$R$ (LLM)} & \textbf{$R$ (Full)} & \textbf{RMSE} \\
\midrule
Relevance & BERTScore, Relevance (LLM), Halluc. (LLM) & 0.673 & 0.610 & 0.585 & 0.688 & 0.944 \\
Recall & BERTScore, Relevance (LLM), METEOR & 0.502 & 0.377 & 0.281 & 0.528 & 1.070 \\
Redundancy & ROUGE-1, Relevance (LLM), METEOR & 0.502 & 0.377 & 0.281 & 0.528 & 1.070 \\
Coherence & ROUGE-2, METEOR, BERTScore & 0.437 & 0.406 & 0.224 & 0.471 & 1.072 \\
Grammar & ROUGE-L, METEOR, BERTScore & 0.420 & 0.377 & 0.247 & 0.463 & 0.826 \\
Hallucinations & ROUGE-2, ROUGE-L, BERTScore & 0.306 & 0.226 & 0.217 & 0.355 & 1.750 \\
\bottomrule
\end{tabular}
\end{table}

The results of the correlation analysis show the expected increase in metrics for a model that incorporated a full stack of predictors. It is noteworthy that among the most significant predictors for expert scores, the BERTScore and Relevance criteria were the most frequent. According to the correlation analysis, the automatic evaluation metrics and the LLM judges' scores demonstrated the highest correlation with the expert scores in terms of Relevance and Recall.

\section{Discussion}\label{sec:discussion}

Relatively high expert scores were observed for ``Grammar, Terminology'' and ``Coherence, Structure'' across most of the tested models. All selected LLMs demonstrated the ability to generate summaries that met the high standards set by the established parameters. The low correlation between the expert scores and LLM-as-a-judge may be explained by both minor differences in the quality of summaries for these criteria and limitations of expert review. Evaluations derived from expert opinion are inherently subjective, being grounded in individual professional experience.

The LLM judges effectively identify significant shortcomings in the summaries; however, they frequently conflate or misapply the evaluation criteria. Interestingly, BERTScore demonstrated a significantly higher correlation with expert opinions compared to LLM-as-a-judge.

The key reason may lie in the fundamental difference in the approaches to scoring. An LLM judge, even when asked to abide by strict criteria, remains a ``generalist'' whose decisions are influenced by the vast array of data on which it was trained. While models simulate reasoning, they suffer from high entropy and are biased toward superficial text attributes. Furthermore, our experiment used a quantized version of Mistral Small 24B, which in itself reduces calculation accuracy and decreases sensitivity to subtle semantic differences.

At the same time, the BERTScore architecture is designed to process text fragments using semantic frame embeddings, maintaining accuracy even with large context lengths. Thus, BERTScore is a ``special tool'' optimized for one specific problem: determining the extent to which the semantics of a single text corresponds to a larger semantic space optimized for the medical domain.

It should also be noted that neither the proposed automatic evaluation metrics nor LLM judges are capable of reliably detecting hallucinations, which occur in LLMs even at zero temperature for summary generation. For medical tasks, this aspect is critical. Errors in the summaries of medical history, laboratory data, and drug therapy can mislead physicians, distort the clinical presentation, and ultimately harm patients.

The analysis revealed several factual inaccuracies in LLM-generated summaries. For example, instances of unacceptable fabrication were observed. In one case, a summary produced by the Gemma-2-27b-it model included data on patient alcohol abuse---information absent from the source text. In an effort to construct a coherent medical narrative, the model apparently substituted the missing detail with a statistically probable inference derived from patterns in its training data.

Finally, LLMs demonstrated notable weaknesses in numerical analysis. Consequently, the LLM-as-a-judge failed to detect a significant error where the reported body temperature was inaccurately stated as $39.0\,^\circ\mathrm{C}$ instead of the correct $36.6\,^\circ\mathrm{C}$. Such a numerical substitution is clinically important.

Modern literature addresses the problem of automated approaches to assessing the credibility of generated texts. Methods for automated detection of hallucinations based on pre-trained models have been proposed~\cite{urlana2025,park2025}. Promising directions are general methods that are not tied to particular domains and do not require training on labeled data. For example, Farquhar et al.~\cite{farquhar2024} propose a method based on semantic entropy, which provides a partial solution for identifying errors associated with confabulations. Verga et al.~\cite{verga2024} propose an approach that employs a diverse panel of models, demonstrating high agreement with expert evaluations.

\section{Conclusions}\label{sec:conclusions}

Despite significant progress in the field of LLMs, full automation of the assessment of medical text summarization performance remains an open question, primarily because neither automatic evaluation metrics nor LLM judges are capable of reliably identifying factual errors and semantic distortions (hallucinations).

At the same time, high expert scores indicate the significant potential of LLMs for medical text summarization tasks. Even open-source architectures with small models have demonstrated the ability to adequately systematize and generalize complex and heterogeneous medical data, indicating great potential for further development.

It is worth noting that the development of an autonomous quality assessment tool does not require full correlation with expert opinion, which appears methodologically challenging due to subjectivity of the latter. Implementation of tailored hallucination detection methods and the use of LLMs with a larger number of parameters or more highly specialized LLMs trained primarily on medical data will improve the reliability, validity, and accuracy.

Nevertheless, even at the current stage, the combination of LLMs and automatic evaluation metrics can provide the basis for medical decision support systems to improve experts' performance by drawing attention to identified deficiencies in medical summaries.

\section*{Limitations and Future Research}

This study employed only low-parameter, local, open-access general-purpose models without fine-tuning, as they present the most straightforward and transparent option for practical deployment. Such models can be easily integrated into existing medical information infrastructures: implementation requires installing a local inference server using Ollama, after which summarization may be obtained via a simple HTTP request.

The research was conducted on a limited dataset using a constrained set of LLMs for both summarization and evaluation. Optimization of model and prompt, followed by validation on a larger dataset, could help identify the most effective approaches for the autonomous assessment of medical summaries.

Another direction for future research involves integrating retrieval-augmented generation (RAG) into the LLM judge architecture. RAG can enhance evaluation objectivity and grounding by retrieving relevant context from an external knowledge base. Further analysis of the influence of hyperparameters, temperature, and model size on the ability to generate and detect hallucinations is also needed.

%
\begin{credits}
\subsubsection{\ackname}
This paper was prepared by a team of authors within the framework of a scientific and practical project in the field of medicine (No.~EGISU:~125051305989-8), ``A promising automated workplace of a radiologist based on generative artificial intelligence.'' The authors express their gratitude to MIREA---Russian Technological University for assistance in conducting the study.

\subsubsection{\discintname}
The authors declare no competing interests to declare that are relevant to the content of this article.
\end{credits}
%
% ---- Bibliography ----
%
\bibliographystyle{splncs04}
\begin{thebibliography}{20}
\bibitem{roustan2025}
Roustan, D., Bastardot, F.: The clinicians' guide to large language models: A general perspective with a focus on hallucinations. Interact. J. Med. Res. \textbf{14}, e59823 (2025)

\bibitem{vanveen2024}
Van Veen, D., et al.: Adapted large language models can outperform medical experts in clinical text summarization. Nat. Med. \textbf{30}, 1134--1142 (2024)

\bibitem{nazi2024}
Nazi, Z.A., Peng, W.: Large Language Models in Healthcare and Medical Domain: A Review. Informatics \textbf{11}(1), 57 (2024)

\bibitem{vasilev2025scoping}
Vasilev, Y.A., et al.: Application of large language models in radiological diagnostics: A scoping review. Digit. Diagn. \textbf{6}, 268--285 (2025)

\bibitem{kaiser2025}
Kaiser, P., et al.: The interaction of structured data using openEHR and large language models for clinical decision support in prostate cancer. World J. Urol. \textbf{43}, 67 (2025)

\bibitem{nakaura2024}
Nakaura, T., et al.: The impact of large language models on radiology: A guide for radiologists on the latest innovations in AI. Jpn. J. Radiol. \textbf{42}, 685--696 (2024)

\bibitem{filali2025}
Filali Ansary, R., Lechien, J.R.: Clinical decision support using large language models in otolaryngology: A systematic review. Eur. Arch. Otorhinolaryngol. \textbf{282}, 4325--4334 (2025)

\bibitem{niel2025}
Niel, O., et al.: Performance evaluation of large language models in pediatric nephrology clinical decision support. Pediatr. Nephrol. \textbf{40}, 3211--3218 (2025)

\bibitem{lee2023}
Lee, P., Bubeck, S., Petro, J.: Benefits, Limits, and Risks of GPT-4 as an AI Chatbot for Medicine. N. Engl. J. Med. \textbf{388}, 1233--1239 (2023)

\bibitem{liu2025}
Liu, F., et al.: Application of large language models in medicine. Nat. Rev. Bioeng. \textbf{3}, 445--464 (2025)

\bibitem{lin2004}
Lin, C.Y.: ROUGE: A package for automatic evaluation of summaries. In: Text Summarization Branches Out, pp. 74--81. ACL, Barcelona (2004)

\bibitem{papineni2002}
Papineni, K., et al.: BLEU: A method for automatic evaluation of machine translation. In: Proc. 40th Annual Meeting of the ACL, Philadelphia (2002)

\bibitem{lavie2007}
Lavie, A., Agarwal, A.: METEOR: An automatic metric for MT evaluation with high levels of correlation with human judgments. In: Proc. Second Workshop on Statistical Machine Translation, Prague (2007)

\bibitem{zhang2019}
Zhang, T., et al.: BERTScore: Evaluating text generation with BERT. arXiv preprint arXiv:1904.09675 (2019)

\bibitem{tang2023}
Tang, L., et al.: Evaluating large language models on medical evidence summarization. npj Digit. Med. \textbf{6}, 158 (2023)

\bibitem{zheng2023}
Zheng, L., et al.: Judging LLM-as-a-judge with MT-bench and chatbot arena. arXiv preprint arXiv:2306.05685 (2023)

\bibitem{chen2024}
Chen, Y., et al.: Comparative analysis of open-source language models in summarizing medical text data. In: Proc. 2024 IEEE Int. Conf. on Digital Health (ICDH), Shenzhen (2024)

\bibitem{croxford2025}
Croxford, E., et al.: Automating evaluation of AI text generation in healthcare with a large language model (LLM)-as-a-judge. medRxiv (2025)

\bibitem{vanschaik2024}
Van Schaik, T.A., Pugh, B.: A field guide to automatic evaluation of LLM-generated summaries. In: Proc. SIGIR '24, Washington, DC (2024)

\bibitem{gu2024}
Gu, J., et al.: A survey on LLM-as-a-Judge. arXiv preprint arXiv:2411.15594 (2024)

\bibitem{urlana2025}
Urlana, A., et al.: HalluCounter: Reference-free LLM hallucination detection in the wild. arXiv preprint arXiv:2501.05376 (2025)

\bibitem{park2025}
Park, S., et al.: Steer LLM latents for hallucination detection. arXiv preprint arXiv:2501.03229 (2025)

\bibitem{farquhar2024}
Farquhar, S., et al.: Detecting hallucinations in large language models using semantic entropy. Nature \textbf{630}, 625--630 (2024)

\bibitem{verga2024}
Verga, P., et al.: Replacing judges with juries: Evaluating LLM generations with a panel of diverse models. arXiv preprint arXiv:2410.14669 (2024)
\end{thebibliography}
\end{document}
